%% kapa: exam=1 round=30 year=2019 period=1 form=A subject=economics subject_ko=경제학원론 seq=26 number=66
%% source: Q-Net 감정평가사 기출문제 (cst00309) artlSeq=5205205 file=extracted/1008/5205205/2019년 제30회 감정평가사 1차시험 문제지/2019년 제30회 감정평가사 1차시험 문제지/2019년도 제30회 감정평가사 1차 1교시 A형.hwp
%% conversion: kordoc 4.24.51 -> tools/gen1.py (auto); review=proofread-sonnet; flags=none
\begin{question}{66}
피셔(I. Fisher)의 기간 간 선택(intertemporal choice)모형에서 최적소비선택에 관한 설명으로 옳은 것을 모두 고른 것은? (단, 기간은 현재와 미래이며, 현재소비와 미래소비는 모두 정상재이다. 무차별곡선은 우하향하며 원점에 대하여 볼록한 곡선이다.)
\begin{framedtext}
ㄱ. 실질이자율이 상승하면, 현재 대부자인 소비자는 미래소비를 증가시킨다. \par
ㄴ. 실질이자율이 하락하면, 현재 대부자인 소비자는 현재저축을 감소시킨다. \par
ㄷ. 실질이자율이 상승하면, 현재 차입자인 소비자는 현재소비를 감소시킨다. \par
ㄹ. 미래소득이 증가하여도 현재 차입제약에 구속된(binding) 소비자의 현재소비는 변하지 않는다.
\end{framedtext}
\begin{choices}
  \item ㄱ, ㄴ
  \item ㄴ, ㄷ
  \item ㄷ, ㄹ
  \item ㄱ, ㄷ, ㄹ
  \item ㄴ, ㄷ, ㄹ
\end{choices}
\end{question}
\answer{4}
